\cotefile{12}
\scanpage{2}
\mspage{1}
\uline{Esquisse d'une théorie des motifs}.

1. Revue des théories cohomologiques habituelles des schémas projectifs lisses.

\margin{$*$}L'idée de motif, \inserted{qui doit remplacer la cohomologie entière (la théorie des modules de Lef. sur $\unsure{\mathrm Z|\mathbb Z}$ !).} Réserves : a) Isomotifs, b) Motifs $\frac12$ simples c) Base ponctuelle.

2. Définition de la catégorie $M^+(k)$ des motifs effectifs sur $k$.
C'est une catégorie abélienne \inserted{``$\frac12$ simple''} avec $\otimes$. Foncteurs $T_{\mathrm{DR}}^{\deleted{\unsure{*}}}$, $T_{\mathrm{Hdg}}$ (si $k$ de car nulle) $T_B$ (si $k=\mathbb C$) $T^\ell$ (si $\ell$ \deleted{\unsure{premier}} $\inserted{\ne\operatorname{car} k}$)\deleted{\illegible{2}} les foncteurs sont multiplicatifs. Tous ces foncteurs sont fidèles.

3. Graduation dans la catégorie $M^+(k)$. \inserted{Le produit $\otimes$ compatible avec graduation.} Les foncteurs $h^i(X)$. $M^{\unsure{i0|i,0|+}}(X)$.

4. L'équivalence $M^{+1}(k)^\circ \simeq \mathrm{Isab}(k)$ ($X\leadsto\underline{\mathrm{Pic}}^{00}_{X}$, $A\leadsto h^1(\hat A)$).

\deleted{4. Les motifs $Q(-1)$, et la catégorie $M(k)$ des motifs.}

\margin{$*$}Les motifs de poids $\geq1$ comme généralisation des VA. Le caractère ``discret'' est conservé (s'exprimant également par le fait que si $k\unsure{\longrightarrow|\lhook\joinrel\longrightarrow} k'$ est une extension de corps algébriquement clos, alors $M^{\unsure{e}}(k)\unsure{-|\longrightarrow}M^e(k')$ est pleinement fidèle).

\margin{$*$}Caractère continu\ldots, Exemples d'énoncé suggéré par la théorie des \deleted{\unsure{schémas}} abéliens : caractérisation des motifs \uline{lisses} sur $S$ \deleted{\unsure{régulier (point|réguliers (point} \unsure{exclu)}}\inserted{normal}.

5. Le motif $Q(\unsure{+1|-1})$ (anti-Tate). La catégorie $M(k)$ des motifs (isomotifs semi-simples sur $k$ !). Maintenant il y a des \uline{Hom} internes, en plus \unsure{des|du} $\otimes$ et de la graduation.

Si $C$ est une catégorie pseudo-abélienne avec $\otimes$, alors les foncteurs compatibles avec $\otimes$ de $M(k)$ dans $C$ correspondent aux foncteurs $T:V(k)^\circ\longrightarrow C$, satisfaisant Künneth, et tels que \deleted{\unsure{l'objet}} $T(Q(-1))=\operatorname{Im} T(\pi)$ ($\pi$ le projecteur dans $h(\unsure{\mathrm P|\mathbb P}^1_k)$ \deleted{\unsure{degré de}}, d'image $h^2(\mathrm P^1_k)=Q(-1)$) \uline{soit} \uline{inversible}. On montre qu'un tel $T$ commute automatiquement à la formation du \uline{Hom}. La catégorie des motifs $M(k)$ est donc caractérisée, avec sa structure $\otimes$, \deleted{et l'objet $Q(\unsure{-1})$} à équivalence près par propriété universelle.
% archive sheet number: 1

\scanpage{3}
\mspage{2}
6. \uline{Niveau d'un motif} \inserted{Résulte de la présence de $M^+$ dans $M$.} Un motif $M$ est dit \uline{de niveau} $\leq i$ s'il est de la forme $M'(n)$, où $M'\in\operatorname{Ob}M^+(k)$ est à composantes homogènes de degrés $\leq i$.
\margin{Exemple : si $i>n=\dim X$, alors $h^i(X)$ est de \uline{niveau} $\leq 2n-i<i$}
On dit que $M$ est purement de niveau $i$ s'il est de niveau $\leq i$, et si aucun sous-motif $\ne0$ n'est de niveau $\leq i-1$. \inserted{Cet $i$ est unique et $\geq0$.} Tout motif \inserted{$M$} se décompose \inserted{de façon unique} en somme de motifs purement de niveau $i$
\[
M=\coprod_{i\in\mathrm Z}M_i
\]
donc s'il est de \deleted{degré} poids $n$, \deleted{il s'écrit} \inserted{les $i$ \deleted{\unsure{tels}} que $M_i\ne0$ ont la $=$ parité que $n$, et}
\[
M=\coprod_i M'_i\left(\frac{\unsure{n}-i}{2}\right)\qquad M'_i\in\mathord{\tikz[baseline=(c.base)]{\node[draw,ellipse,inner sep=2pt](c){$M^{+i}_{\mathrm{pur}}(k)$};}}
\]
où $M'_i$ est \deleted{\unsure{purement de poids}} \inserted{\uline{effectif}} pur de poids et niveau $i$.

\deleted{\unsure{Il est facile}} Les motifs de \deleted{poids}\inserted{niveau} $\leq i$ sont ceux qui, modulo twist de Tate, ``proviennent'' de variétés $X\in V(k)$ de dimension $\leq i$. Ceci résulte du fait que $h(X)$ est de niveau $\leq\dim X$, et que \deleted{\unsure{si $i$|Addition}} si $i\leq\dim X$, $h^i(X)\unsure{\longleftrightarrow|\lhook\joinrel\longrightarrow}h^i(Y)$, où $Y$ est une section linéaire de dim $i$ (\unsure{en vertu|au sens|sens} du ``th. de Lefschetz'').
\margin{\uline{Ex} Si $i\geq n=\dim X$ $h^i(X)$ est de niveau $\leq2n-i$}
Autre interprétation en termes de cohomologie $l$-adique \deleted{(en supposant $k$ alg. clos, ce qui est loisible, car la notion de niveau, pureté du \unsure{m} niveau, est invariante par extension du corps de base).} On regarde dans $H^i(\bar X,Q_\ell)$ la ``filtration géométrique'' définie par
\[
\mathrm{Filt}^{\unsure{j|J}}H^i(\bar X,Q_\ell)=
\underset{U,\;U\text{ ouvert tel que }\operatorname{codim}(X-U)\geq \unsure{j|d}}{\underline{\lim}}\,
\operatorname{Ker}\bigl(H^i(\bar X,Q_\ell)\to H^i(U,Q_\ell)\bigr)
\]
% archive sheet number: 2

\scanpage{4}
\mspage{3}
Alors on a, \deleted{si $X$ purement de dim $n$ et $i\leq n$ (resp. $i\geq n$)}
\[
\begin{aligned}
\mathrm{Filt}^{\unsure{j|d}} H^i(\bar X,Q_\ell)&=\text{partie de }T^\ell(h^i(X))\ \text{\deleted{provenant}\inserted{image} du}\\
&\quad\text{plus grand sous-motif \inserted{de}}\\
&\quad\text{niveau }\leq i-2j\ \text{\deleted{(resp. $2n-i\unsure{-j|\geq j}$)} \deleted{de} de }h^i(X)
\end{aligned}
\]
Cette caractérisation \inserted{du niveau des motifs} est tout au moins correcte si on admet la résolution des singularités, \deleted{\unsure{sur $k$}} \inserted{des variétés de dim $\leq\dim X=n$ sur $k$.} Comme conséquence de la résolution et des conjectures standard, on trouve donc que les dimensions des $\mathrm{Filt}^j H^i(\bar X,Q_\ell)$ sont indépendantes de $\ell$. (-- Bien mieux, ils correspondent à des motifs, et les quotients facteurs successifs sont \deleted{\unsure{et}} purs de poids $i-2j$ resp. $2n-i-j$).

\inserted{Parler des inv. birationnels\ldots}
% slanted
\margin{\uuline{Conjectures de Tate et de Hodge pourraient passer ici}}
7. \uline{Structure des} $\otimes$-\uline{catégories}. \deleted{\unsure{abéliennes|standard}}

\deleted{$\otimes$ catégorie.} Catégorie abélienne avec opération $\otimes$ ``ass'', ``comm'' et ``unitaire''. On suppose de plus l'existence de \uline{Hom} internes, au sens du $\otimes$.

On dit qu'elle est ``spéciale'' si
\[
\check M\otimes N\xrightarrow{\sim}\underline{\mathrm{Hom}}(M,N)
\]
Supposons de plus que $C$ soit $K$-linéaire ($K$ un corps) (\unsure{i.e.} on se donne $K\longrightarrow\operatorname{End}(\mathrm{id}_{C})$), et que les Hom soient de dim finie sur $k$ : \deleted{\unsure{on est}} On dira qu'on a une $\otimes$-catégorie sur $k$ (pour simplifier).

Une telle catégorie est \unsure{dite|dit} ``connexe'' si $\operatorname{End}(\mathbf{1}_C)\simeq k$ (\unsure{par|bien} l'identité).
% archive sheet number: 3

\scanpage{5}
\mspage{4}
Exemple : Si $G$ est un groupe algébrique sur $k$, la catégorie $\mathrm{Modf}(G)$ des vectoriels sur lesquels $K$ opère. Item avec groupes proalgébriques (p. ex. profinis).

Cet exemple \deleted{est} essentiellement universel, comme on va voir.

$\otimes$-foncteurs. Foncteur\unsure{s} fibre\unsure{s} \inserted{sur $K$, -- sur $K'$, \ldots.} Considérons un couple $(C,F)$ \deleted{\unsure{$s/F$|formé}} d'une $\otimes$-catégorie sur $K$, avec foncteur fibre \unsure{ntr.|etc.|rat.} /$K$. À cette situation \deleted{\unsure{est}} on peut associer un groupe \deleted{alg} affine $G$ sur $K$ (ou pro-groupe pro-affine sur $K$), et une équivalence $(C,F)\simeq(\mathrm{Modf}(G),\text{oubli})$.
% slanted
\margin{$G(A)=\operatorname{Aut}F_A$\quad $F_A(M)=F(M)\otimes_K A$ \deleted{\unsure{considéré}} comme foncteur fibre à valeurs dans $\mathrm{Mod}(A)$}
De façon précise, associant à chaque $G$ le couple $\bigl((\mathrm{Mod}(\unsure{f|G}),G),\text{oubli}\bigr)$, on trouve un foncteur
\diagram{\begin{tikzcd}[ampersand replacement=\&,column sep=large]
\text{groupe affine sur }K \arrow[r] \& \begin{gathered}\otimes\ \text{catégorie avec foncteur}\\\text{fibre donné, avec }\otimes\\\text{-foncteurs mod isom,}\\\text{compatibles avec }F\end{gathered}
\end{tikzcd}}
et ce foncteur est pleinement fidèle.
\[
\begin{aligned}
C\ \text{``de type fini''}&\iff G\text{ de type fini}\\
C\ \tfrac12\text{-simple}&\underset{\inserted{\text{si }K\text{ de car. nulle}}}{\iff}\deleted{\unsure{G\ \text{fini réductif}}}\ G\text{ proréductif}.\\
&\phantom{\iff}\deleted{\unsure{\text{fibre algébrique}}}
\end{aligned}
\]
\uline{NB} Possibilité d'une définition des autres op. tensorielles telles $\wedge^{\unsure{i}}$, $\operatorname{Sym}^i$ (au fait, on voit que c'est indépendant de l'existence et du choix de $F$)

Si on a un foncteur fibre, tout autre foncteur fibre $F'$ s'obtient en ``tordant'' avec un pro-torseur $P$ sous $G$, $P=(P_i)$ :
\[
F'=P\times^G F
\]
% slanted
\margin{\uuline{à prouver} \unsure{Revient|Reviens|Remarque} : affirmer que si $C$ de type fini, deux foncteurs fibres deviennent isom. sur ext. finie}
% archive sheet number: 4

\scanpage{6}
\mspage{5}
Alors $G'=P\times^G G=\operatorname{ad}(P)$

En fait, il y a des cas intéressants où \deleted{si} on peut trouver aucun foncteur fibre. Mais si $C$ est de type fini, alors on \inserted{doit pouvoir} montrer qu'on peut trouver un foncteur fibre sur une extension finie de $K$. \deleted{On \unsure{réussirait|construit} ainsi à décrire les} \inserted{Quand il en est ainsi, \deleted{on \unsure{doit avoir|démontre}}} admettant le résultat ``d'unicité'' de foncteurs fibres, \deleted{on doit} arriver à décrire $(C,F)$ intrinsèquement en termes d'une gerbe sur un site convenable\ldots

8. \uline{Cas des catégories des motifs}. \deleted{\uline{Lien avec le groupe fondamental profini}.}

$S$ schéma régulier \inserted{\unsure{rédiger|intègre|connexe}}. On définit $M^+(S)$, $M(S)$ en paraphrasant ce qui est fait sur un corps, en travaillant (pour $X$ projectif lisse sur $S$, avec $C(X)\overset{\mathrm{dfn}}=C(X_\eta)$ ($\eta$ pt générique de $S$). Ainsi, $M(S)\unsure{\longrightarrow|-}M(\eta)$ est pleinement fidèle, permettant d'interpréter $M(S)$ comme ss-catégorie pleine de $M(\eta)$. De plus, pour tt morphisme $f:S'\longrightarrow S$ de schémas réguliers, on déduit $f^*:M(S)\longrightarrow M(S')$ un $\otimes$-foncteur correspondant (c'est ici que sert l'hypothèse de \uline{régularité} sur les schémas\ldots.).
% slanted
\margin{à expliciter\ldots.}

Tout point \inserted{géom} \uline{complexe} $\xi$ de $S$ définit un foncteur fibre ``de Betti'' \deleted{$\unsure{F^\xi_{\mathrm{Betti}}=F_\xi}$} $F^{\mathrm{Betti}}_\xi=F^B_\xi:M(S)\unsure{-|\longrightarrow}\mathrm{Modf}(\unsure{Q|\mathbb Q})$, d'où \unsure{un|une} groupe proalgébrique $G$ sur $Q$, proréductif
% archive sheet number: 5

\scanpage{7}
\mspage{6}
en fait, \deleted{et} décrivant $M(S)$. Pratiquement, on n'a guère qu'à considérer \deleted{dans} \deleted{\unsure{sur|une}} chaque question une ss-catégorie de type fini de $M(S)$, soit $C$. Elle est alors décrite par un \deleted{\unsure{groupe}} groupe algébrique \inserted{réductif} $G$ sur $Q$. Question intéressante : comment sont reliés ces groupes pour des $\xi$ différents (ils se déduisent tous les uns des autres par twist\ldots). Déjà intéressant si $S$ est le spectre d'une ext. finie $k$ de $Q$\ldots. !

Si $S$ est de car. $p>0$, on ne pourra en général trouver de foncteur fibre sur $Q$ (contre-exemple de Serre, lié aux courbes elliptiques de Hasse nul). Mais on aura un foncteur fibre sur $Q_\ell$ pour tt pt géom de car. $\ne\ell$, d'où un foncteur fibre sur une extension finie \inserted{$K$} de $Q$ (si on se restreint à une $C\subset M(S)$ stable par \deleted{\unsure{$\otimes$, fibre}} $\otimes,\vee$, et à \uline{\unsure{engendrt|engendrée}} fini\ldots). Donc on trouve un groupe algébrique réductif $G_K$ sur $K$, qui permet de représenter la catégorie des $K'$-motifs pour tt corps de base $K'$ qui est une extension de $K$.

\margin{////}
\inserted{Structure $\varepsilon,\unsure{\eta|D|y}$, \unsure{$\varepsilon\eta$\textup{ gr}$=\lambda^{?}$|$\varepsilon y(\lambda)=\lambda^2$}, structure $\unsure{\eta_1,\eta_2|\gamma_1,\gamma_2|y,i_1,i_2}$}
\deleted{\unsure{pour la compréhension du|pour corps de base} $G$, \unsure{$T^*$|$T_{\mathrm{Hdg}}$}}
\unsure{9)} \uline{Relations avec} \uline{le} $\pi_1$ \uline{\deleted{profini}\inserted{Hdg}}, : \inserted{A) $G/G^0\unsure{-|\longrightarrow}\pi^{1}$.} Soit $F$ un foncteur fibre à coeff$^{\mathrm{ts}}$ dans $K$ sur $M(C)$, $\xi$ un pt géométrique. On dit que $F$ et $\xi$ sont associés si on a un isom
% archive sheet number: 6

\scanpage{8}
\mspage{7}
\[
F\mid M^{+0}(S)\simeq F_\xi
\]
où $F_\xi$ est le foncteur évident sur $M^{+0}(S)$ : fibre en $\xi$. (NB. $S$ étant connexe, ces foncteurs sont isom entre eux, donc la \uline{condition} d'association ne dépend pas de $\xi$. Mais il faut \unsure{le|la} considérer plutôt comme une \uline{donnée} supplémentaire, comme \uline{donnée} de l'isom en question).

\uline{Ex 1} Si $\operatorname{source}\xi=\operatorname{Spec}(\mathbb C)$, \deleted{alors} $F=F^{\mathrm{Betti}}_\xi$, alors on a une telle donnée. \inserted{\uline{Ex. 2} $T^\ell_\xi$}

\uline{Ex} \deleted{\unsure{2}}3 Si $S=\operatorname{Spec}k$, $k$ corps de car. nulle (p. ex. $k=Q$) et si $F$ est le foncteur fibre \deleted{sur}\inserted{à coefficients dans} $k$ $T_{\mathrm{DR}}$ \deleted{alors}, ou $T_{\mathrm{Hdg}}$, alors il n'est pas en général associé aux pts géométriques.

\uline{Ex} 4 Si \deleted{$k$ est} $S$ est le spectre d'un corps alg clos, alors tt $F$ est associé à $\xi$, du moins si $k$ est dénombrable \inserted{i.e. de degré de transc. absolu dénombrable}, donc $M(k)$ à engendrement dénombrable\ldots .

\deleted{Alors on trouve un isom canonique
\[
\pi_1(S,\xi)\simeq\varprojlim G_\alpha/G_\alpha^0
\]
\inserted{Les $G_\alpha/G_\alpha^0$ sont des groupes constants sur $K$ \ldots. :}
Cela résulte du \deleted{fait que} $\otimes$-foncteur canonique
\[
\mathrm{Mod}\unsure{\text{ adm}|\text{ disc}|\text{ dis}|\text{ cont}}(\pi_1)\simeq M^{+0}(S)\longrightarrow M(S)
\]
qui définit
\[
G\longrightarrow\pi_1
\]}
% archive sheet number: 7

\scanpage{9}
\mspage{8}
Alors on trouve un $\otimes$-foncteur \inserted{pleinement} fidèle
\diagram{\begin{tikzcd}[ampersand replacement=\&]
\pi_1\text{-}Q\text{-modules admissibles de rang fini sur }Q \arrow[r] \& \text{motifs sur }S
\end{tikzcd}}
qui est décrit par un \inserted{épi}morphisme
\[
G=\varprojlim G_\alpha\unsure{\twoheadrightarrow|\longrightarrow}\pi_1(S,\xi)
\]
et se factorise donc en un \uline{épimorphisme} canonique
\[
\boxed{\varprojlim_{\alpha}G_\alpha/G_\alpha^0\longrightarrow\pi_1(S,\xi)}\tag*{(*)}
\]
Soit $\pi'_1$ le premier mb. Alors \deleted{les $\pi'_1$-} \inserted{(supposant $K=Q$} pour simplifier) les $\pi'_1$-$Q$-modules correspondent à la sous-catégorie de $M(S)$ formée des motifs \inserted{$M$} tels que la sous-\inserted{$\otimes$-}catégorie engendrée par $M$ soit de ``longueur finie''. Utilisant conjecture de Tate (ou de Hodge si $S$ de car. générique nulle) on \inserted{peut} \unsure{voir|voit} que ce sont les motifs \deleted{\unsure{de}} de niveau \inserted{et poids} zéro, \inserted{i.e. les motifs effectifs de poids zéro,} i.e. que $(*)$ \uline{est un isomorphisme}. Les \deleted{résultats}\inserted{réflexions} précédents peuvent se \deleted{\unsure{décrire}}\inserted{développer} aussi (utilisant la notion de motif sur $S$ à coeff. dans $K$) si on ne suppose pas $K=Q$.

\deleted{\uline{NB} La bijectivité de \unsure{ce signe}}
% letterhead
% printed upside down, lower right
INSTITUT\\
DES HAUTES ETUDES SCIENTIFIQUES\\
91 - BURES-SUR-YVETTE\\
928 48-53

\uline{NB} Soit $S'\longrightarrow{}S$ un revêtement étale \inserted{\unsure{surj}}\deleted{de} $S$, \inserted{$S'$ connexe,} et supposons que $F$ provienne d'
% archive sheet number: 8

\scanpage{10}
\mspage{9}
un \deleted{revêtement} $\otimes$-foncteur \deleted{fibre} sur $S'$. Alors, comme en théorie \unsure{du groupe fondamental ordinaire|des groupes fondamentaux ordinaires} $G'\longrightarrow G$ est un monomorphisme, dont l'image est l'\deleted{\unsure{image|noyau}} image inverse \inserted{dans $G$} du sous groupe ouvert de $\pi_1(S,\xi)$ qui correspond à $S'/S$. On en conclut que la bijectivité de $(*)$ \deleted{équivaut (pour \unsure{tt revêtement} $S$\unsure{-|=})} (pour \deleted{\unsure{tout|tt}} $S$, disons) équivaut au fait que pour $S=\operatorname{Spec} k$, $k$ alg. clos, le \deleted{\unsure{$G$}} groupe $G$ est \uline{connexe}.

10) \uline{Relations avec le} $\pi_1$ \uline{profini} : \uline{les hom} $\pi_1\to G(Q_\ell)$.

Supposons \deleted{\unsure{encore}} que $K\underset{i}{\hookrightarrow}Q_\ell$ ; \deleted{et que}\inserted{on dit que} $F$ \unsure{est|soit} \deleted{associé au pt géom. $\xi$.} \uline{compatible avec le foncteur de Tate} $T^\ell_\xi$ (pour le plongement donné de $K$ dans $Q_\ell$) si $F\otimes_K Q_\ell\xrightarrow{\sim}T^\ell_\xi$ (donnée de compatibilité avec foncteur de Tate).

\uline{Exemple} 1 \ $K=Q_\ell$, $F=T^{\unsure{1|\ell}}_\xi$,

\uline{Exemple} 2 \ $K=Q$, $F=T_{\mathrm{Betti}}^\xi$, où $\xi$ un pt complexe.

Alors \deleted{\unsure{repr.}} l'homomorphisme canonique
\[
\pi_1(S,\xi)\to\operatorname{Aut}T^\ell_\xi
\]
\deleted{définit un} joint à $\operatorname{Aut}T^\ell_\xi\simeq\operatorname{Aut}F\otimes_{K}Q_\ell\simeq G(Q_\ell)$.
définit un homomorphisme
\[
\boxed{\rho_\ell:\pi_1(S,\xi)\longrightarrow G(Q_\ell)}
\]
% archive sheet number: 9

\scanpage{11}
% letterhead
INSTITUT\\
DES HAUTES ETUDES SCIENTIFIQUES\\
91 - BURES-SUR-YVETTE\\
928 48-53
% archive sheet number: 10

\scanpage{12}
\mspage{10}
Noter que si $F$ est compatible avec $T^\ell_\xi$, il est à fortiori \uline{associé à} $\xi$, et on trouve
\[
G\deleted{\unsure{(Q_\ell)}}\longrightarrow\pi_1(S,\xi)
\]
d'où
\[
\unsure{\eta|\gamma}_\ell:G(Q_\ell)\longrightarrow\pi_1(S,\xi).
\]
On voit tout de suite que cet homomorphisme est inverse à gauche de $\unsure{\rho|\varphi}_\ell$
\[
\boxed{\eta_\ell\rho_\ell=\mathrm{id}_{\pi_1(S,\xi)}}
\]
Si on s'intéresse à une sous-catégorie $C_\alpha$ de $C$, correspondant à un $G_\alpha^{\unsure{0}}$, alors $G_\alpha/G_\alpha^0$ \deleted{correspond à}\inserted{définit} un quotient $\pi_\alpha$ de $\pi_1(S,\xi)$, et on a
\diagram{\begin{tikzcd}[ampersand replacement=\&,column sep=small,row sep=small]
\& \& G_\alpha/G_\alpha^0\arrow[dr,dashed] \& \\
\pi_1(S,\xi)\arrow[r,"\rho_\ell"]\arrow[rrr,bend right=27,"\text{\uline{hom can.}} "'] \& G_\alpha(Q_\ell)\arrow[ur,dashed]\arrow[rr] \& \& \pi_\alpha
\end{tikzcd}\quad(commutatif).}
La conjecture de Tate \inserted{(avec action semi-simple de Galois !)} équivaut à dire que, si $S$ est localisé d'un schéma de type fini sur $\mathrm Z$, alors $G_\alpha\bigl(\unsure{\to|-}Q_\ell\bigr)$ est l'enveloppe algébrique de $\operatorname{Im}(\pi_1(S,\xi)\unsure{-|\to}G_{\unsure{\ell|\alpha}}(Q_\ell))$.
% slanted
\margin{\uline{NB} La conjecture de Tate sous la forme : les $\frac12$-invariants dans $H^{2i}(\bar X,Q_\ell(i))$ sont les comb. à coeff. dans $Q_\ell$ des classes algébriques -- implique (avec conjectures standard) la $\frac12$ simplicité dans Tate\ldots}
\margin{\deleted{NB \unsure{il s'ensuit|$G_\alpha(Q_\ell)$} \illegible{3} \unsure{que l'action} \unsure{sur} $T_\ell(M)$, \unsure{algébriquement} \unsure{engendrée}}}

\uline{Cas} $S=\operatorname{Spec} k$, $k$ \uline{corps fini}.

Soit $C\subset M(k)$, et $F$ un ``foncteur fibre sur $C$ à coeff dans $K$. Comme on a un élément canonique \deleted{de}
\[
\mathrm{Frob}_k\in\operatorname{Aut}\underline{\mathrm{Id}}_{M(k)}\longrightarrow\operatorname{Aut}\mathrm{id}_C
\left[\begin{gathered}\simeq\unsure{Z|\mathcal Z|\mathfrak Z}(Q)\\
\text{où }Z\text{ est le centre de }G\\
\deleted{\unsure{C/F}}
\end{gathered}\right]
\]
on en déduit un élément canonique de $G(K)=\operatorname{Aut}F$, et même de $Z(Q)$ (où $Z$ est le centre de $G$)
\[
\boxed{\unsure{\xi|*}\ \mathrm{Frob}_{k\unsure{,|/}F}\in Z(Q)\subset G(K)}.
\]
\margin{N.B. $Z$ est en fait défini sur $Q$ !}
% archive sheet number: 11

\scanpage{13}
\mspage{11}
Si $K\subset Q_\ell$ et si $F$ est compatible avec $T^\ell_\xi$, alors l'image de $\mathrm{Frob}_{k,F}\in \unsure{Z|\mathcal Z|\mathfrak Z}(K)\unsure{\subset|=|\to}Z(Q_\ell)\subset G(Q_\ell)$ \deleted{n'est autre} n'est autre que l'image du \deleted{\unsure{produit sur|produit par}} générateur
\[
\mathrm{Frob}_k\in\pi_1(k,\xi)\simeq\widehat{\mathrm Z}
\]
par l'hom
\[
\unsure{\rho|\varphi}_\ell:\widehat{\mathrm Z}\longrightarrow G(Q_\ell)
\]
La \uline{conjecture de Tate} pour le corps fini $k$ équivaut à dire que $G$ est l'enveloppe algébrique du \unsure{ss-groupe|groupe} engendré par $\mathrm{Frob}_{k,F}\in G(K)$. Elle implique donc que $G$ est commutatif. [Comme il ne dépend que de \unsure{twists|twist}, on voit qu'il est en tous cas défini sur $Q$, canoniquement, et évidemment on aura en fait, alors
\[
\mathrm{Frob}_{k,F}\in\deleted{Z(Q)\subset}G(Q)\quad ].
\]
\uline{Conséquence de conjecture de Tate pour} $k$.

$G$ est extension \deleted{\inserted{(canoniquement trivialisée)}} de $\widehat{\mathrm Z}$ par un tore resp. protore.

\uline{Conséquence} : Soit à nouveau $S$ quelconque, mais muni d'un point fermé $s$ tel que $k(s)=k$ soit fini. \deleted{Soit} $C\subset M(S)$, \deleted{et soit $C'$ une sous-catégorie de}\inserted{muni} d'un foncteur fibre $F$, à coeff. dans un corps $K\supset Q$,
% archive sheet number: 12

\scanpage{14}
\mspage{12}
d'où un $G$ schéma en groupes sur $K$. Soit $I_G$ le ``schéma des invariants'', isomorphe si $\unsure{C|G}$ connexe \inserted{à} $T/W$ ($T$ tore maximal \unsure{?|de $G$}, $W$ son groupe de Weyl).

Alors on trouve un élément $f_s$ de $I_G(Q)$. (NB. $I_G$ est en fait défini sur $Q$). On trouve ainsi une famille d'éléments $(f_s)_{s\in S^0}$, où $S^0$ désigne l'ens des pts de $S$ tels que $k(s)$ soit fini : famille des frobenius. Compatibilité avec les hom
\[
G\deleted{\unsure{(K)}}\longrightarrow\pi_1(S,\xi)
\]
et les hom
\[
\pi_1(S,\xi)\longrightarrow G(Q_\ell).
\]
% archive sheet number: 13

\scanpage{15}
\mspage{13}
12) Structures supplémentaires sur le groupe de Galois motivique : \inserted{structures associées aux} graduations, filtrations, au motif de Tate.

\deleted{La catégorie} On choisit une sous-\inserted{$\otimes$-}catégorie graduée $C$ de $M(S)$, avec un foncteur fibre $F$, d'où groupe de Galois motivique $G$

a) Une graduation de type $\unsure{\pi|\Pi|M}$ sur le foncteur fibre \inserted{à coeff dans $K$} compatible avec $\otimes$, est équivalent à la donnée d'\unsure{une|un} hom
\[
T=D_K(\pi)\longrightarrow G
\]
La graduation provient d'une graduation fonctorielle de $M\in\inserted{\operatorname{Ob}C}$ lui-\unsure{$=$|même} \unsure{sss|ssi} l'hom $T\longrightarrow G$ est central. Par suite :

$\underline{1^\circ})$ La graduation \inserted{canonique} de type $\mathrm Z$ des objets de $C$ définit un hom canonique
\[
\unsure{\eta|\gamma|y}:\mathrm{G}_{m\,K}\longrightarrow G
\]
$2^\circ)$ Si $F=T_{\mathrm{Hdg}}$, la bigraduation de ce foncteur correspond à un hom
\[
\deleted{D_K(\mathrm Z^2)=\mathrm G_m^2\longrightarrow G,\quad\text{donc à deux h}}
\]
\[
i:\mathrm G_{m\,K}^2\longrightarrow G
\]
donc à deux hom commutant l'un à l'autre
\[
i_1,i_2:\mathrm G_m\unsure{\longrightarrow|-}G
\]
Comme la graduation totale associée est celle de $1^\circ)$, on trouve
\[
i_1(\lambda)i_2(\lambda)=\eta(\lambda)
\]
(ce qui entraîne déjà que $i_1,i_2$ commutent\ldots).
% archive sheet number: 14

\scanpage{16}
\mspage{14}
$3^\circ)$ Supposons $S=\operatorname{Spec}(\mathbb C)$, prenons $F=T_{\mathrm{Betti}}$ \deleted{\unsure{ou} $T_{\mathrm{Hdg}}$} \deleted{\unsure{foncteur Hdg|Bett / Hdg}.} Alors on \deleted{peut prendre} $K=Q$. D'autre part $F\otimes_Q\mathbb C\simeq T_{\mathrm{Hdg}}\deleted{\unsure{\otimes_Q\mathbb C}}$, donc \deleted{\unsure{on|on a}} $G^{\mathrm{Betti}}\otimes_Q\mathbb C\simeq G^{\mathrm{Hdg}}\deleted{\unsure{\otimes_Q\mathbb C}}$. Donc dans $G^{\mathrm{Betti}}\otimes_Q\mathbb C$ on trouve \deleted{\unsure{un}} $i_1,i_2$
\[
i_1,i_2:\mathrm G_m\longrightarrow G^{\mathrm{Betti}}\otimes_Q\mathbb C
\]
tels que
\[
i_1(\lambda)i_2(\lambda)=\unsure{\eta|\gamma|y}(\lambda).
\]
De plus on a
\[
\overline{i_1(\lambda)}=i_2(\lambda)
\]
ce qui exprime que les deux graduations \unsure{entières|partielles|naturelles|relatives} sont conjuguées complexes l'une de l'autre.

b) Il y a lieu de supposer que $\operatorname{Ob}C\ni Q(-1)$. Alors $Q(-1)$ \deleted{\unsure{est}} définit un $G$-module $F(Q(-1))$ de rang 1, d'où un hom
\[
\varepsilon:G\unsure{\longrightarrow|-}\mathrm G_{m\,K}
\]
Comme $Q(-1)$ est de \deleted{degré} poids 2, on aura
\[
\boxed{\varepsilon(\eta(\lambda))=\lambda^2}
\]

c) Supposons que sur le foncteur $F$ il y ait une filtration fonctorielle, compatible avec gr. Comment l'expliciter ? On considère
\[
F'=\mathrm{Gr}(F),
\]
% archive sheet number: 15

\scanpage{17}
\mspage{15}
c'est un autre foncteur fibre, mais \uline{gradué}, donc défini par un \uline{torseur} $P$ \deleted{$\unsure{G}$} sous $G$, et $F$ \unsure{une forme|en termes} de $F'$ \unsure{---} par torseur $P'$ sous $G'$. Supposons que sur \deleted{\unsure{quelque} \unsure{point}} une extension $K'$ de $K$, on puisse trouver un isom $F\sim F'$ compatible avec les filtrations, et induisant l'identité sur les gradués associés.

[ \uline{Ex} \inserted{$S=\operatorname{Spec}k$,} $k$ de car. nulle, $F=T_{\mathrm{DR}}$ avec filtration canonique ; ça marche aussi si $S$ est n'importe quel schéma \inserted{régulier} de car. nulle \ldots ]. Alors il y a un sous-torseur \inserted{canonique} $Q$ de $P$ \deleted{$\unsure{G}$}, de groupe $H$ le sous-groupe de $G$ des autom compatibles avec filtration et induisant l'identité sur le gradué, et de $=$ \deleted{\unsure{pour}} un $Q'$ dans $P'$, \inserted{groupe $H'$.} On notera que $H,H'$ sont unipotents, Ainsi, $F$ se déduit canoniquement du foncteur gradué $F'$ [graduation définie par $\mathrm G_m\xrightarrow[K]{\alpha}G'$] à l'aide d'un torseur $Q'$ sous le groupe \inserted{unipotent} $H(\alpha)$ de $G'$ défini par $\alpha$ comme on pense.

% slanted
\margin{expliciter la construction de $H(\alpha)$}

13) \uline{Relations avec les groupes de \deleted{\unsure{Galois}}\inserted{Mumford-Tate}.}

Les structures de Hodge polarisables forment une $\otimes$-catégorie semi-simple. Toute sous-\deleted{\unsure{$\otimes$}}catégorie de type fini est ``définie'' par un \inserted{groupe alg. $G$ sur $Q$} \textsuperscript{\unsure{b}}.

Soit $(V,V_{\unsure{0|Q}},(V^{pq}))$ une structure de Hodge.

\textsuperscript{\unsure{b}} avec \unsure{$i'_1,i'_2$|$i_1,i_2$} : $\mathrm G_{m\,\unsure{\mathbb C|\mathrm C}}\unsure{-|\longrightarrow}G_{\mathbb C}$ satisfaisant $i_1(\lambda)i_2(\lambda)=i(\lambda)$ central, $i_2(\lambda)=\overline{i_1(\lambda)}$.

On a \unsure{réciproquement}
% archive sheet number: 16

